% Incorporación focal de dependencias anteriores; no sustituye al núcleo común.
% Procedencia detallada: technical/CIERRE_DEPENDENCIAS_ANTERIORES_VI.md.
% Los valores de prueba se publican después de sus operadores; no son entradas.
\section{Coordonnées engendrées, action et normalisation électronique}
\label{vi:dep:seccion}

La construction des collections utilise des résultats des étapes précédentes :
les caractères régionaux, le registre dodécaphasique, la droite d’action et la
section électronique. Nous réunissons ici les opérations qui les relient et les
preuves nécessaires à leur composition. Leur ordre est causal : les régions et leurs
registres sont produits sur APP par TRIT et TPK ; les évaluations réelles
sont effectuées ensuite ; le choix des unités intervient lors de la réalisation des
grandeurs dimensionnelles.

Les conditions initiales et la sélection régionale ont été développées dans
la section précédant l’arithmétique des histoires. Le domaine du registre
incidentiel, ses dénombrements et la reconstruction du registre sont définis
expressément ci-dessous. La sélection d’une collection concrète de
visites et de traces est une opération supplémentaire à l’évaluation d’un registre :
elle nécessite de fixer ses multiplicités, le transport de frontière et la
couverture des traces admises. Les formules de reconstruction qui sont
démontrées ici permettent d’évaluer et de vérifier ce registre ; son
inversibilité ne sélectionne pas à elle seule la collection incidentielle.

\subsection{Lecteurs régionaux et compatibilité de l’évaluation}
\label{vi:dep:regiones}

La clôture, la propagation et l’auto-échelle sont trois régions d’une même
génération. La région de clôture comporte cinq émissions et une multiplicité
\(432+4\cdot144=1008\). Dans la jauge orientée du catalogue, ses émissions
sont
\[
\begin{array}{c|r}
(501,614,498,169,272,272)&432\\
(810,923,870,169,272,272)&144\\
(810,983,810,169,272,272)&144\\
(870,923,810,169,272,272)&144\\
(870,923,810,418,674,521)&144.
\end{array}
\]
Le mot ternaire commun est \(010211\) ; les représentants de propagation
et d’auto-échelle sont \(201101\) et \(121200\). Le recensement et l’orientation
précèdent les évaluations qui suivent.

Aux cinq positions du lecteur de clôture, la projection uniforme
\(P_5=\mathbf1\mathbf1^{\mathsf T}/5\), avec
\(P_5\lambda=\lambda\) et \(\mathbf1^{\mathsf T}\lambda=1\), détermine
\(\lambda=\mathbf1/5\). Cette normalisation des positions du lecteur
est distincte de la distribution des germes indiquée dans le tableau.
Dans le régime elliptique \(J^2=-I\), la multiplication donne
\[
(I+aJ)(I+bJ)=(1-ab)I+(a+b)J.
\]
Pourvu que \(1-ab\ne0\), la pente composée est
\(a\oplus b=(a+b)/(1-ab)\). Deux duplications de \(1/5\) donnent
\(5/12\) et \(120/119\). Le compensateur positif \(1/q\), \(q>1\),
qui ramène la pente à un satisfait
\[
\frac{120/119-1/q}{1+(120/119)/q}=1,
\qquad q=239.
\]
Par conséquent, le caractère de clôture de ce lecteur est
\begin{equation}
 \Pi_{\rm cl}=4\left(4\arctan\frac15-\arctan\frac1{239}\right).
 \label{vi:dep:lector-clausura}
\end{equation}
La règle de composition, sa compensation et le choix de la branche angulaire
fixent cette valeur ; l’identité angulaire conventionnelle la reconnaît comme \(\pi\).

Le caractère de propagation se réalise dans une collection formelle commutative
\(P_t\), normalisée par \(P_0=1\), \(P'_0=1\) et
\(P_{t+s}=P_tP_s\). Si \(P_t=\sum_{n\ge0}a_nt^n\), la comparaison
du coefficient de \(s^nt\) donne
\[
 (n+1)a_{n+1}=a_n,\qquad a_0=1.
\]
Ainsi, \(a_n=1/n!\) ; la représentation à une unité de paramètre est
\(\Pi_{\rm pr}=\sum_{n\ge0}1/n!\), reconnue comme \(e\).
L’unité et le signe du générateur appartiennent au lecteur : la seule loi
de semi-groupe admettrait d’autres générateurs.

Pour l’auto-échelle, TRIT sélectionne les modes
\((\Sigma,\Pi,\Pi)\). Leurs trois transitions cycliques sont
\(\Sigma\to\Pi\), \(\Pi\to\Pi\) et \(\Pi\to\Sigma\). Leur matrice
d’incidence, dans l’ordre aréal--volumétrique, est
\begin{equation}
 F=\begin{pmatrix}0&1\\1&1\end{pmatrix}.
 \label{vi:dep:incidencia-hojas}
\end{equation}
Le vecteur propre positif normalisé comme \((1,x)^{\mathsf T}\)
satisfait \(x=\lambda\), \(1+x=\lambda x\) ; par conséquent,
\(x^2-x-1=0\). Sa racine positive est la coordonnée d’auto-échelle
\(\varphi\). Cette équation reconnaît le caractère de l’incidence déjà
produite par le calendrier.

\begin{proposition}[Évaluation régionale par intervalles compatibles]
\label{vi:dep:prefijos}
Les trois lecteurs précédents disposent d’intervalles rationnels de largeur
arbitrairement petite. Leurs préfixes décimaux finis sont déterminés par
raffinement et sont compatibles par troncature.
\end{proposition}
\begin{proof}
Pour \(0<x<1\), deux sommes alternées consécutives de
\(\sum_{j\ge0}(-1)^jx^{2j+1}/(2j+1)\) encadrent \(\arctan x\).
Le premier terme omis borne l’erreur. Cela fournit des intervalles
rationnels pour \eqref{vi:dep:lector-clausura}.
Si \(E_n=\sum_{j=0}^n1/j!\), alors, pour \(n\ge1\),
\[
 E_n<\Pi_{\rm pr}<E_n+
 \frac{n+2}{(n+1)(n+1)!}.
\]
La borne résulte de la majoration des quotients successifs de la queue par
\(1/(n+2)\). Enfin, \(F^n\) a pour entrées
\(F_{n-1},F_n,F_{n+1}\), où \(F_0=0,F_1=1\) et
\(F_{n+1}=F_n+F_{n-1}\). L’identité
\(F_{n+1}^2-F_nF_{n+2}=(-1)^n\) montre que deux quotients consécutifs
encadrent \(\varphi\) avec une largeur \(1/(F_nF_{n+1})\).
On peut prendre l’intersection des intervalles successifs pour conserver
l’emboîtement. Après la reconnaissance conventionnelle, l’irrationalité de
\(\pi,e,\varphi\) assure qu’aucune coordonnée ne se trouve sur une frontière
décimale rationnelle. Pour chaque profondeur finie, un intervalle
suffisamment étroit tient dans une unique cellule décimale. Son indice
détermine le préfixe et l’emboîtement démontre leur compatibilité.
\end{proof}

\subsection{Coordinatisation dodécaphasique signée de l’état TPK plein}
\label{vi:dep:estado-pleno}

L’évolution qui produit le registre est formulée sur l’état complet du
Topological Phase Kernel. Notons
\(\mathsf{TPKFullState}\) l’espace des états qui conserve conjointement
position APP, régime TRIT, route, feuillet, orientation, retenue, frontière,
mémoire, registre incidentiel et coordonnées terminales. Dans une coordinatisation canonique,
\[
 X=(X_{\rm APP},X_{\rm TRIT},X_{\rm ruta},X_{\rm mem},
       C_{\rm term},U_{12}^{\rm sgn},X_{\partial},\ldots).
\]
La dynamique de cent huit pas agit par la composition
\begin{equation}
 X_{\rm term}
 =\bigl(\mathcal U_{108}\circ\cdots\circ\mathcal U_1\bigr)(X_{\rm seed}),
 \qquad
 \mathcal U_t=\mathsf{Upd}_t\circ\mathsf{Tra}_t\circ\mathsf{Sel}_t,
 \label{vi:dep:estado-terminal-108}
\end{equation}
où \(\mathsf{Sel}_t\) applique le sélecteur tritique de phase,
\(\mathsf{Tra}_t\) transporte feuillet, route et orientation, et
\(\mathsf{Upd}_t\) incorpore la retenue et la mémoire produites à
l’instant \(t\). Chaque facteur agit donc sur l’état plein et la
composition conserve la généalogie des douze fenêtres nonadiques.

Le registre incidentiel de la section suivante est la coordonnée
\(\mathcal R_{12}(X_{\rm term})\). Son évaluation par les canaux
\(90/120\) définit
\[
 \mathcal E_{108}^{90,120}\circ\mathcal R_{12}:
 \mathsf{TPKFullState}\longrightarrow\mathbb Z^{25},
 \qquad
 X\longmapsto C(X)=(b^{90}(X),b^{120}(X),Q(X)).
\]
La projection dodécaphasique signée est
\[
 \mathcal S_{12}:=\operatorname{pr}_{U_{12}^{\rm sgn}}:
 \mathsf{TPKFullState}\longrightarrow\mathbb Z^{12}.
\]
Sur le sous-réseau compatible produit par la dynamique, les deux coordinatisations
satisfont l’identité opératoire
\begin{equation}
 \mathcal S_{12}
 =\Pi_H\circ\mathcal E_{108}^{90,120}\circ\mathcal R_{12}
 =\Pi_H\circ\operatorname{pr}_{C}.
 \label{vi:dep:identidad-productor-e108}
\end{equation}
Ainsi, la signature dodécaphasique est obtenue par projection et évaluation du même
état qui conserve le registre ; son domaine inclut l’orientation et la mémoire
qui interviennent dans chaque dénombrement. La projection brute
\(\rho_{108}:\mathsf{TPKFullState}\to\mathsf{CT108RawProjection}\)
désigne, en revanche, la coordinatisation qui oublie ces coordonnées après la
production. Cette séparation des types fixe l’ordre causal entre génération,
représentation signée et lecture réduite.

\subsection{Des incidences temporelles au registre dodécaphasique}
\label{vi:dep:registro}

\begin{definition}[Domaine du registre incidentiel orienté]
\label{vi:dep:libro-dominio}
Un registre \(\mathscr L\) contient une collection finie non vide de routes,
leurs visites identifiées et une collection finie de traces. Chaque route
conserve une visite par instant pendant les \(108\) premiers instants,
et les visites ultérieures qu’utilisent ses traces. Une visite conserve
\[
 v=(\gamma,t,i,j,s,\tau,w,\varepsilon_\partial),
\]
où \(\gamma\) identifie la route, \(t\) son instant, \(1\le i,j\le9\)
sa position APP, \(s\in\{\Sigma,\Pi,0\}\) son feuillet,
\(\tau\in\{-1,0,1\}\) son orientation tritique et \(w\) sa multiplicité
entière positive. La coordinatisation incidentielle évalue
\[
 n_v=
 \begin{cases}
 i+j,&s=\Sigma,\\
 ij,&s=\Pi,\\
 9,&s=0.
 \end{cases}
 \qquad
 r_v=1+((n_v-1)\bmod9),\quad q_v=(n_v-r_v)/9.
\]
La lecture \(n_v=9\) fixe la convention de seuil de cette coordinatisation.
Pour \(r_v=9\), l’étiquette \(\varepsilon_\partial\) distingue
couronne et intérieur ; elle ne s’obtient pas uniquement à partir du résidu.

Une trace \(\theta=(v_1,\ldots,v_N)\) parcourt des instants consécutifs d’une
même route. Elle conserve une multiplicité \(u(\theta)\in\mathbb N_{>0}\)
et une unité de longueur réduite
\(\lambda(\theta)\in\mathbb N_{>0}\), avec
\[
 N\lambda(\theta)=120(1+3j(\theta)),\qquad j(\theta)\in\mathbb N_0.
\]
Son canal est
\(\sigma(\theta)=\operatorname{sgn}(\sum_{a=1}^N\tau(v_a))\).
La collection dénombrée dans les canaux \(120\) comprend les traces de
signe non nul. La longueur réduite et le nombre d’instants sont
reliés par \(\lambda\) ; ils ne sont pas identifiés l’un à l’autre.
\end{definition}

La définition spécifie le domaine de l’évaluation, y compris les
données qui doivent être transmises au lecteur. Ses multiplicités et la collection
de traces appartiennent à la sélection dynamique du registre. Une énumération
du domaine ou une vérification de format ne remplace pas cette sélection.

La génération conserve un registre \(\mathscr L\) de visites et de traces.
La coordinatisation temporelle le divise en fenêtres
\(E_m=\{9m-8,\ldots,9m\}\), \(1\le m\le12\),
équivalentes aux positions \(m=\beta+1\) du registre
\(\mathcal R_{12}\). Chaque fenêtre retient sa sous-route,
son orientation, son incrément de frontière et ses incidences locales.
Chaque visite enregistre position APP, feuillet arithmétique, instant, multiplicité
et orientation de frontière. À chaque évaluation \(n_v\) est conservée la
division \(n_v=r_v+9q_v\), avec \(1\le r_v\le9\). Dans la fenêtre
\(m\), le lecteur incidentiel exprime
\begin{align}
 A_m&=\sum_{v\in\mathscr L_m}w(v)
      \mathbf1_{\{1,3,5,7\}}(r_v),\notag\\
 C_m&=\sum_{j\ge0}
       \bigl(\nu^-_{120,m}(j)-\nu^+_{120,m}(j)\bigr),\notag\\
 V_m&=\sum_{v\in\mathscr L_m}w(v)
      \mathbf1_{\{9\}}(r_v)\,\varepsilon_{\partial}(v).
 \label{vi:dep:incidencias}
\end{align}
Ici les traces de chaque fenêtre ont un support fini,
\(\varepsilon_{\partial}=+1\) dans la couronne et \(-1\) à l’intérieur,
et les étiquettes \(\pm,120\) appartiennent au transport conservé.
La multiplicité \(w\) et les traces \(\nu\) sont celles obtenues par ce
transport ; elles ne sont pas reconstruites en choisissant une valeur de \(\alpha\).

En termes de traces identifiées, la seconde somme a pour
expression finie
\[
 C_m=-\sum_{\theta:\,t(v_1)\in E_m}u(\theta)\sigma(\theta).
\]
En effet, regrouper par \(j(\theta)\) et par signe produit
\(\nu^-_{120,m}(j)-\nu^+_{120,m}(j)\). Cette égalité spécifie
l’opération effectivement appliquée à chaque registre ; elle conserve
séparés le canal, la multiplicité et la longueur.

\begin{proposition}[Mise à jour et concaténation du registre]
\label{vi:dep:libro-actualizacion}
La collection de visites et de traces admises par le transport étant fixée, leurs
dénombrements sont mis à jour par des incréments entiers. Lorsqu’on ajoute une
visite \(v\) de la fenêtre \(m\), l’incrément de \((A_m,C_m,V_m)\) est
\[
 \left(w(v)\mathbf1_{\{1,3,5,7\}}(r_v),\,0,\,
 w(v)\mathbf1_{\{9\}}(r_v)\varepsilon_\partial(v)\right).
\]
Lorsqu’une trace \(\theta\) dont l’origine appartient à \(E_m\) est achevée,
l’incrément est \((0,-u(\theta)\sigma(\theta),0)\).
Pour deux segments consécutifs, le dénombrement de leur concaténation est
la somme des dénombrements des visites et des traces contenues dans chaque
segment, plus celui des traces qui franchissent leur jonction.
\end{proposition}
\begin{proof}
Chaque somme de \eqref{vi:dep:incidencias} est indexée par des identifiants
de visites ou de traces ; ajouter un identifiant ajoute donc
exactement le terme écrit. Pour la concaténation, les visites
se répartissent entre les deux segments. Les traces se répartissent en
celles contenues dans le premier, celles contenues dans le second et celles
qui traversent la jonction. L’additivité des sommes entières prouve
la formule. Conserver les identifiants et les prolongements
en attente évite de compter deux fois une trace ou d’en perdre une qui
n’est pas encore terminée. La réduction décimale est appliquée après le
dénombrement ; les quotients de cette réduction restent dans la mémoire.
\end{proof}

La réduction décimale de ces trois dénombrements produit
\begin{equation}
 z_m=100[A_m]_{10}+10[C_m]_{10}+[V_m]_{10},
 \qquad 0\le z_m<1000.
 \label{vi:dep:bloque-incidencial}
\end{equation}
Soient \((D_3z)_i=z_i-z_{i+3}\),
\((D_4z)_i=z_i-z_{i+4}\), avec des indices dans \(\ZZ/12\ZZ\), et
\(Q(z)=\sum_i z_i\). Le lecteur conserve
\(\mathcal D(z)=(D_3z,D_4z,Q(z))\). La réduction décimale oublie des
multiples de dix de chaque dénombrement, tandis que le registre original conserve
les visites et les transports qui les ont engendrés.
La lecture des événements de code pair sur cette même histoire,
sa décomposition \(12=1+5+6\) et son inversion entière sont
développées dans la proposition \ref{vi:prop:eventos}. Ce vecteur
d’événements ne remplace pas le présent registre d’incidences :
les domaines et les coordonnées conservées restent distincts.

\begin{proposition}[Conjugaison incidentielle et réduction décimale]
\label{vi:dep:conjugacion-incidencial}
Soit \(\varrho(m)=13-m\). Considérons une conjugaison du registre
qui envoie les visites et les traces de la fenêtre \(m\) sur celles de
\(\varrho(m)\) par des bijections qui conservent leurs multiplicités.
Supposons qu’elle conserve la classe arithmétique et l’étiquette de frontière
des visites, et échange les canaux positif et négatif des
traces. Alors les dénombrements de \eqref{vi:dep:incidencias} satisfont
\begin{equation}
 (A_m^\dagger,C_m^\dagger,V_m^\dagger)
 =(A_{\varrho(m)},-C_{\varrho(m)},V_{\varrho(m)}).
 \label{vi:dep:conjugacion-recuentos}
\end{equation}
En écrivant \(u_m=[C_m]_{10}\), leurs blocs décimaux vérifient
\begin{equation}
 z_m^\dagger=z_{\varrho(m)}
       +100\mathbf1_{\{u_{\varrho(m)}\ne0\}}-20u_{\varrho(m)}.
 \label{vi:dep:conjugacion-bloque}
\end{equation}
\end{proposition}
\begin{proof}
Les bijections conservent chaque terme de \(A_m\) et \(V_m\).
L’échange des canaux permute \(\nu^-\) et \(\nu^+\),
ce qui change le signe de \(C_m\). On obtient
\eqref{vi:dep:conjugacion-recuentos}. Si \(C=10q+u\),
\(0\le u<10\), la division euclidienne de \(-C\) donne
\[
 [-C]_{10}=10\mathbf1_{\{u\ne0\}}-u,
 \qquad
 \left\lfloor\frac{-C}{10}\right\rfloor
 =-q-\mathbf1_{\{u\ne0\}}.
\]
Dans \eqref{vi:dep:bloque-incidencial}, seul le chiffre
intermédiaire change. Sa différence est
\(10([-C]_{10}-[C]_{10})=100\mathbf1_{\{u\ne0\}}-20u\),
ce qui démontre \eqref{vi:dep:conjugacion-bloque}. La seconde
identité conserve simultanément le quotient entier de la réduction.
\end{proof}

La proposition précise l’action sur le registre avant de transporter
sa représentation décimale. Son application à l’inversion d’une route
conserve également la convention de lecture des extrémités. Soit
\(\gamma=(x_0,\ldots,x_n)\) une route enrichie et soit \(f\) une
observable entière évaluée avant chaque avancée. La route inverse
\(\gamma^{-1}=(x_n,\ldots,x_0)\) vérifie
\begin{equation}
 \sum_{k=0}^{n-1}f(x_{n-k})
 =\sum_{k=0}^{n-1}f(x_k)+f(x_n)-f(x_0).
 \label{vi:dep:conjugacion-frontera}
\end{equation}
En effet, le membre de gauche somme les indices \(1,\ldots,n\) ;
le premier terme de celui de droite somme \(0,\ldots,n-1\). Leur différence
est exactement le terme de frontière indiqué. L’identité s’applique
à chaque observable incidentielle avec ses données conservées. Les extrémités
sont comptabilisées dans les dénombrements entiers avant réduction modulo dix ;
sur une route fermée le terme s’annule, tandis que dans une fenêtre
ouverte il peut modifier son bloc. Les traces qui traversent une jonction
sont conservées au moyen de la partition de la
proposition~\ref{vi:dep:libro-actualizacion}.

Ainsi les dénombrements sont reliés à la conjugaison des routes de la
section~\ref{vi:sec:conjugacion}, en maintenant déclarée l’action
concrète sur chaque incidence. La négation d’un canal et l’inversion
d’une route sont des opérations sur le registre ; leur réduction décimale est la
représentation qui vient d’être calculée.

\begin{proposition}[Image entière et reconstruction du registre]
\label{vi:dep:imagen-integral}
Soient \(b,c\in\ZZ^{12}\), \(q\in\ZZ\), avec des indices de zéro à onze.
Définissons
\[
 w_i=c_i-b_{i+1},\qquad P_0=0,\qquad
 P_i=\sum_{j=0}^{i-1}w_j,\qquad T_w=\sum_{i=0}^{11}P_i.
\]
Il existe un unique \(k\in\ZZ^{12}\) tel que
\(D_3k=b\), \(D_4k=c\), \(Q(k)=q\), si et seulement si
\begin{equation}
 b_i=w_i+w_{i+1}+w_{i+2},\qquad
 \sum_iw_i=0,\qquad q+T_w\equiv0\pmod{12}.
 \label{vi:dep:condiciones-imagen}
\end{equation}
Dans ce cas,
\begin{equation}
 k_i=\frac{q+T_w}{12}-P_i.
 \label{vi:dep:inversa-incidencial}
\end{equation}
\end{proposition}
\begin{proof}
Si \(k\) existe, alors
\(c_i-b_{i+1}=k_i-k_{i+1}\). La somme cyclique de ces différences
est nulle, leur somme de longueur trois est \(b_i\), et
\(k_i=k_0-P_i\). En sommant les douze coordonnées,
\(q=12k_0-T_w\), ce qui prouve les conditions nécessaires et la formule.
Réciproquement, les conditions rendent la formule entière et
cyclique. Ses différences de longueur trois sont \(b_i\), et celles de
longueur quatre sont \(w_i+b_{i+1}=c_i\). Sa somme est \(q\).
La différence entre deux solutions a toutes ses différences
consécutives nulles et sa somme nulle ; elle est donc nulle.
\end{proof}

La reconstruction est une opération sur l’image typée, non une
sélection rétrospective de ses arguments. En particulier,
\(\mathcal D(z)=\mathcal D(z')\) implique \(z=z'\).
Comme l’expression d’un entier de \([0,999]\) en trois chiffres est unique,
deux triplets de dénombrements expriment le même registre si et seulement s’ils coïncident
composante par composante modulo dix.

Le changement de représentation signé utilise le bloc
\begin{equation}
 H_4=\begin{pmatrix}
 1&1&1&1\\1&1&-1&-1\\
 1&-1&1&-1\\1&-1&-1&1
 \end{pmatrix},\qquad H_4^2=4I_4.
 \label{vi:dep:hadamard}
\end{equation}
L’opérateur \(H_{12}\) applique \(H_4\) aux trois orbites ordonnées
\((1,4,7,10)\), \((2,5,8,11)\), \((3,6,9,12)\).
Sur le sous-réseau
\(\mathcal L_H=\{U\in\ZZ^{12}:H_{12}U\equiv0\pmod4\}\),
les applications
\begin{equation}
 U=H_{12}K,\qquad K=\tfrac14H_{12}U
 \label{vi:dep:K-reversible}
\end{equation}
sont inverses. La congruence exprime exactement le caractère entier de la
seconde application, et \(H_{12}^2=4I\) prouve les deux compositions.

Le registre signé est également calculé directement à partir des coordonnées
transversales, avant de reconstruire \(K\). Avec des indices de zéro à onze,
soient, pour \(r=0,1,2\),
\[
 S_r=\sum_{j=0}^3c_{r+3j},\quad
 a_0=\frac{q+2S_0+S_1}{3},\quad
 a_1=a_0-S_0,\quad a_2=a_1-S_1,
 \qquad d_{r,j}=b_{r+3j}.
\]
L’application harmonique exprime par orbites
\begin{equation}
 \Pi_H^{(r)}(b,c,q)=
 (a_r,d_{r,1}-d_{r,3},d_{r,0}+d_{r,2},d_{r,0}-d_{r,2}).
 \label{vi:dep:PiH}
\end{equation}
Pour prouver la formule sur l’image entière, soit \(k\) son
antécédent. Si \(s_r=\sum_jk_{r+3j}\), alors
\(S_r=s_r-s_{r+1}\) et \(q=s_0+s_1+s_2\), avec des indices d’orbite
modulo trois. La résolution donne \(s_r=a_r\). Pour les quatre coordonnées
\((x_0,x_1,x_2,x_3)\) d’une orbite,
\(d_j=x_j-x_{j+1}\). Leurs trois combinaisons dans
\eqref{vi:dep:PiH} sont respectivement
\(x_0+x_1-x_2-x_3\), \(x_0-x_1+x_2-x_3\) et
\(x_0-x_1-x_2+x_3\). Avec \(a_r=\sum_jx_j\), elles sont
exactement \(H_4(x_0,x_1,x_2,x_3)^{\mathsf T}\).
La composition effective est ainsi
\(\mathscr L\to(b,c,q)\to\Pi_H(b,c,q)=U\to H_{12}U/4=K\).
L’antécédent utilisé dans cette démonstration vérifie l’identité ;
ce n’est pas un argument de l’application \(\Pi_H\).

Pour l’état terminal de \eqref{vi:dep:estado-terminal-108}, la coordonnée
produite avant la représentation signée est
\begin{align}
 b^{90}={}&(-495,-116,-684,108,601,-90,
              -173,-88,313,560,-397,461),\notag\\
 b^{120}={}&(-425,-281,-481,671,-255,30,
               475,-543,680,251,6,-128),\notag\\
 Q_{\rm TPK}={}&6263.
 \label{vi:dep:canales-terminales}
\end{align}
En substituant ces vingt-cinq champs dans \eqref{vi:dep:PiH}, les trois
orbites signées sont
\begin{align*}
 U_{\rm term}^{(1)}&=(2378,-452,-668,-322),\\
 U_{\rm term}^{(2)}&=(1406,998,-204,-28),\\
 U_{\rm term}^{(3)}&=(2479,-551,-371,-997).
\end{align*}
L’entrelacement par colonnes produit
\begin{equation}
 U_{\rm term}=\mathcal S_{12}(X_{\rm term})
 =(2378,1406,2479,-452,998,-551,-668,-204,-371,-322,-28,-997).
 \label{vi:dep:U-terminal}
\end{equation}
La transformée orbitale \(H_{12}/4\) est définie sur ce sous-réseau par
la divisibilité exacte de chaque bloc de quatre. Par conséquent, la chaîne
\begin{equation}
 X_{\rm term}\xrightarrow{\mathcal R_{12}}\mathscr L(X_{\rm term})
 \xrightarrow{\mathcal E_{108}^{90,120}}C_{\rm term}
 \xrightarrow{\Pi_H}U_{\rm term}
 \xrightarrow{H_{12}/4}K
 \label{vi:dep:cadena-terminal-completa}
\end{equation}
est une composition typée allant de l’état plein à son sceau périodique.
Les propositions~\ref{vi:dep:libro-actualizacion} et
\ref{vi:dep:imagen-integral} fournissent, respectivement, la mise à jour
finie du registre et l’inverse entière de la dernière représentation.

Le registre engendré fournit le témoin ordonné
\begin{equation}
 K=(234,543,140,729,659,824,621,058,914,794,146,601).
 \label{vi:dep:K-testigo}
\end{equation}
Le zéro initial de \(058\) fait partie de la lecture par blocs.
Pour le vérifier au moyen de la proposition~\ref{vi:dep:imagen-integral}, ses différences
et sa somme sont
\[
\begin{split}
b={}&(-495,-116,-684,108,601,-90,-173,-88,313,560,-397,461),\\
c={}&(-425,-281,-481,671,-255,30,475,-543,680,251,6,-128),\\
q={}&6263.
\end{split}
\]
Dans les trois orbites précédentes, le registre signé est
\[
 \begin{aligned}
 U^{(1)}&=(2378,-452,-668,-322),\\
 U^{(2)}&=(1406,998,-204,-28),\\
 U^{(3)}&=(2479,-551,-371,-997).
 \end{aligned}
\]
Appliquer \(H_4/4\) à ces trois vecteurs reproduit
\eqref{vi:dep:K-testigo}. Ces identités vérifient la représentation des
incidences ; leur genèse reste celle du registre
\eqref{vi:dep:incidencias}.

\subsection{\texorpdfstring{Représentation périodique de \(K\) et normalisation de \(\alpha\)}{Représentation périodique de K et normalisation d’alpha}}
\label{vi:dep:alpha}

Fixons \(B=1000\), l’origine de phase et les douze positions. Écrivons
\[
 N_K=\sum_{j=1}^{12}K_jB^{12-j},\qquad
 \kappa=\frac{N_K}{B^{12}-1}.
\]
C’est l’évaluation périodique exacte du registre. La multiplication
par \(B^{12}-1\) et la division euclidienne successive récupèrent ses douze
blocs, y compris tout zéro initial. Pour
\(\rho K=(K_2,\ldots,K_{12},K_1)\), la même concaténation donne
\(\kappa(\rho K)=B\kappa(K)-K_1\). Par conséquent, la réversibilité
du registre conserve une origine de phase spécifiée.

\begin{proposition}[Période minimale de la représentation rationnelle]
\label{vi:dep:periodo-minimo-k}
Le prolongement périodique du registre \eqref{vi:dep:K-testigo}
a une période minimale de \(12\) blocs en base \(1000\) et de
\(36\) chiffres en base \(10\).
\end{proposition}
\begin{proof}
Les douze composantes de \(K\) sont distinctes. Toute translation
cyclique non nulle inférieure à douze identifierait deux composantes si elle était
une période. Par conséquent, la période minimale des blocs est douze.

La concaténation de leurs blocs de trois chiffres, y compris \(058\),
produit une période décimale de longueur \(36\). Ce développement est
canonique : il n’est pas ultimement constant au chiffre \(9\).
Soit \(d\) sa période décimale minimale. Alors \(d\mid36\),
car les décalages qui préservent un mot périodique forment
un sous-groupe cyclique. En regroupant les chiffres trois par trois, la période
minimale des blocs est le plus petit \(k>0\) tel que \(d\mid3k\), c’est-
à-dire \(d/\gcd(d,3)\). Aucun diviseur propre de \(36\)
ne satisfait \(d/\gcd(d,3)=12\). Il s’ensuit que \(d=36\).
\end{proof}

Cette période appartient à la coordonnée rationnelle du registre, avec
l’origine de phase fixée. La composition suivante conserve en outre les
coordonnées régionales et les quotients de normalisation de sa sortie.

Soient \(P_j,E_j,\Phi_j\) les blocs fractionnaires des trois
caractères régionaux déjà construits, et prolongeons \(K_j\) avec une période de
douze. La composition des registres est
\begin{equation}
 Z_j=P_j+E_j-\Phi_j-K_j,
 \qquad
 A_j=Z_j+c_{j+1}-Bc_j.
 \label{vi:dep:normalizacion-alpha}
\end{equation}
Le quotient entier \(c_j\) transporte la continuation vers la position
observée. Dans une fenêtre finie, on fixe d’abord \(c_{N+1}\), puis
\[
 c_j=\left\lfloor\frac{Z_j+c_{j+1}}B\right\rfloor,
 \qquad A_j=Z_j+c_{j+1}-Bc_j.
\]
La division euclidienne détermine de manière unique \(0\le A_j<B\).
Comme \(-1998\le Z_j\le1998\), l’ensemble
\(\{-2,-1,0,1,2\}\) est stable pour ces quotients.

\begin{theorem}[Composition exacte à profondeur illimitée]
\label{vi:dep:alpha-completa}
L’évaluation complète de la normalisation, dans la convention décimale
canonique, est
\begin{equation}
 \alpha_A=\pi+e-\varphi-4-\kappa.
 \label{vi:dep:alpha-A}
\end{equation}
Il existe un unique développement canonique \((A_j)\), avec des blocs entre zéro
et \(B-1\), et une suite bornée de quotients entiers qui satisfont
\eqref{vi:dep:normalizacion-alpha} et \(c_1=0\).
\end{theorem}
\begin{proof}
Les parties entières régionales sont \(3,2,1\). Par convergence absolue,
\[
 y:=\sum_{j\ge1}Z_jB^{-j}
   =(\pi-3)+(e-2)-(\varphi-1)-\kappa.
\]
Les premiers blocs sont \(141,718,618,234\), de sorte que
\(Z_1=7\). La somme de la queue a une valeur absolue au plus
égale à \(1998\sum_{j\ge2}B^{-j}=0.002\) ; en particulier \(0<y<1\).
Prenons le développement canonique \(y=\sum A_jB^{-j}\), en choisissant le
développement terminal et non celui ultimement constant \(B-1\) si une
frontière se présente. Posons
\[
 c_j=\sum_{r\ge0}(Z_{j+r}-A_{j+r})B^{-r-1}.
\]
L’égalité des valeurs complètes montre que chaque \(c_j\) est
l’opposé d’une somme finie d’entiers pondérés par des puissances de
\(B\) ; il est donc entier, et \(c_1=0\). Séparer le premier terme
donne \(Bc_j=Z_j-A_j+c_{j+1}\). Les bornes des queues rendent la
suite bornée. Plus précisément, la queue de \(Z\) est dans \([-2,2]\)
et la queue canonique de \(A\) dans \([0,1)\), de sorte que
\(-3<c_j\le2\). Étant entier, \(c_j\) appartient à l’ensemble
\(\{-2,-1,0,1,2\}\) indiqué.
Réciproquement, la somme de la récurrence jusqu’à \(N\) donne
\[
 \sum_{j=1}^N A_jB^{-j}
 =\sum_{j=1}^N Z_jB^{-j}-c_1+c_{N+1}B^{-N}.
\]
Le caractère borné annule le dernier terme à la limite. L’unicité du
développement canonique de \(y\) détermine \(A\), et la formule des queues
détermine alors tous les \(c_j\).
\end{proof}

L’évaluation finie \(\kappa_{36}=N_K/B^{12}\) diffère de
\(\kappa\) de \(B^{-12}\kappa\). En outre, imposer
\(c_{13}=0\) dans une fenêtre de douze blocs ne remplace pas le quotient
de la continuation complète. Dans le registre exprimé, la continuation
donne \(c_{13}=-1\) : le douzième bloc est \(662\), tandis que la
fenêtre tronquée avec un quotient terminal nul donne \(663\). L’évaluation
complète commence par
\[
 \alpha_A=0.007297352569283800997285105472380662999\ldots.
\]
La période de \(K\) n’impose pas de période à \((A_j)\), car dans
\eqref{vi:dep:normalizacion-alpha} interviennent également les trois
registres régionaux et les quotients entiers.

Dans les formules suivantes, \(\alpha\) désigne cette sortie positionnelle
complète \(\alpha_A\). La voie analytique des vacances possède son opérateur
propre ; un polynôme tronqué d’ordre neuf et la racine de ce polynôme
ne sont pas identifiés ici à l’opérateur complet. Les compositions de
masses qui suivent ne nécessitent pas de promouvoir cette approximation en
une égalité entre les deux voies infinies.

\subsection{Construction tridimensionnelle du préfacteur électronique}
\label{vi:dep:prefactor}

Soit \(\mathcal A_3=\{1,\ldots,9\}^3\), muni de la mesure uniforme, et
\(\rho_9(n)=1+((n-1)\bmod9)\). Sur ce domaine agissent
\[
 \Sigma(x)=\rho_9(x_1+x_2+x_3),\qquad
 \Pi(x)=\rho_9(x_1x_2x_3).
\]
Dans \(\ell^2(\mathcal A_3)\), soient \(P\) et \(Q\) les moyennes
conditionnelles sur les fibres de \(\Sigma\) et \(\Pi\),
respectivement. Chacune est idempotente et autoadjointe : dans chaque
fibre, elle remplace un vecteur par sa composante constante et conserve
orthogonalement cette composante. Ce sont donc des projecteurs orthogonaux.

Les fibres additives ont pour cardinal \(81\). Les multiplicatives ont
pour cardinaux
\[
m=(36,36,108,36,36,108,36,36,297).
\]
L’incidence conjointe
\(N_{ab}=\#\{x:\Sigma(x)=a,\Pi(x)=b\}\) est \(9N_0\), où
\begin{equation}
 N_0=\begin{pmatrix}
0&1&1&0&1&1&0&1&4\\1&0&1&1&0&1&1&0&4\\
1&0&2&0&1&2&0&0&3\\0&1&1&0&1&1&0&1&4\\
1&0&1&1&0&1&1&0&4\\0&0&2&1&0&2&0&1&3\\
0&1&1&0&1&1&0&1&4\\1&0&1&1&0&1&1&0&4\\
0&1&2&0&0&2&1&0&3
\end{pmatrix}.
 \label{vi:dep:incidencia-cubica}
\end{equation}
Cette matrice est obtenue en dénombrant les \(729\) triplets. Elle est également
reproduite, sans arrondi, avec
\(N_{ab}=\sum_{i,j,k=1}^9
\mathbf1_{\{a\}}(\rho_9(i+j+k))
\mathbf1_{\{b\}}(\rho_9(ijk))\).
Dans les bases normalisées de fonctions constantes sur chaque fibre,
l’entrelacement a pour matrice \(C_{ab}=N_{ab}/\sqrt{81m_b}\).

\begin{proposition}[Norme de la non-commutativité des deux feuillets]
\label{vi:dep:norma-PQ}
Les projecteurs précédents satisfont
\begin{equation}
 \|[P,Q]\|=\frac{\sqrt3}{4}.
 \label{vi:dep:prefactor-valor}
\end{equation}
\end{proposition}
\begin{proof}
Dans l’espace des neuf fibres additives, \(PQP\) a pour matrice
\(M=CC^{\mathsf T}\). La multiplication des matrices de
\eqref{vi:dep:incidencia-cubica} donne la valeur propre un sur le vecteur
constant, \(1/4\) sur \((1,-1,0,1,-1,0,1,-1,0)\), et
\(7/132\) sur \((1,1,-2,1,1,-2,1,1,-2)\).
Sur les coordonnées \(\{3,6,9\}\) de somme nulle, il agit par
\(1/18\), avec une multiplicité de deux. Sur chacun des ensembles
\(\{1,4,7\}\) et \(\{2,5,8\}\), les vecteurs de somme nulle
appartiennent au noyau ; ils donnent quatre dimensions. Ces sous-espaces
orthogonaux totalisent les neuf dimensions.

Pour un vecteur propre unitaire \(u\in\operatorname{ran}P\) de
\(PQP\) de valeur propre \(0<\lambda<1\), le vecteur
\(v=(Qu-\lambda u)/\sqrt{\lambda(1-\lambda)}\) est unitaire et
orthogonal à \(\operatorname{ran}P\). Dans la base \((u,v)\),
\[
 P=\begin{pmatrix}1&0\\0&0\end{pmatrix},\quad
 Q=\begin{pmatrix}\lambda&\sqrt{\lambda(1-\lambda)}\\
 \sqrt{\lambda(1-\lambda)}&1-\lambda\end{pmatrix}.
\]
Le commutateur a pour norme \(\sqrt{\lambda(1-\lambda)}\).
Les blocs correspondant à \(\lambda=0,1\) commutent. Parmi les
valeurs propres calculées, le maximum est atteint en \(\lambda=1/4\),
car \(\lambda(1-\lambda)\) est croissante sur \([0,1/2]\).
Sa racine carrée est \(\sqrt3/4\).
\end{proof}

Dans le bloc extrémal, \(S=2P-I\) et
\(J=4[P,Q]/\sqrt3\) satisfont
\[
 S^2=I,\qquad J^2=-I,\qquad SJ=-JS.
\]
Les quatre matrices \(I,S,J,SJ\) sont linéairement indépendantes,
de sorte qu’elles réalisent \(\mathrm{Cl}_{1,1}\simeq M_2(\RR)\).
Les éléments \(n_\pm=(S\pm J)/2\) vérifient
\(n_\pm^2=0\) et \(n_+n_-+n_-n_+=I\), par développement direct.
Le préfacteur de l’équation électronique provient ainsi de deux évaluations
d’APP en dimension trois et de leur commutateur, avant toute masse
de référence.

\subsection{Lecture de frontière et composition électronique basale}
\label{vi:dep:basal}

L’incidence \eqref{vi:dep:incidencia-hojas} possède un vecteur propre
positif \(r=(1,\varphi)^{\mathsf T}\). La normalisation
\(p_{ij}=F_{ij}r_j/(\varphi r_i)\) produit
\[
 P_{\rm av}=\begin{pmatrix}0&1\\\varphi^{-2}&\varphi^{-1}\end{pmatrix},
 \qquad \mu=\frac{(1,\varphi^2)}{1+\varphi^2}.
\]
L’identité \(\varphi^{-2}+\varphi^{-1}=1\) prouve que les lignes
ont pour somme un. Comme \(F\) est symétrique, les poids proportionnels à
\(r_i^2\) satisfont l’équilibre détaillé ; ainsi \(\mu P_{\rm av}=\mu\).
L’incidence est irréductible et comporte une autorétention, de sorte que
cette distribution stationnaire est unique. La fréquence chronologique
\((1/3,2/3)\) compte les trois instants du calendrier ; \(\mu\)
pondère ses trajectoires admissibles de mémoire un.

L’opérateur qui conserve les feuillets, fixe la normalisation volumétrique
à un et marque l’aréale par le caractère de clôture est
\(D_\pi=\operatorname{diag}(\pi,1)\). Par récurrence,
\[
 F^n=\begin{pmatrix}F_{n-1}&F_n\\F_n&F_{n+1}\end{pmatrix}.
\]
Le quotient des retours diagonaux marqués est
\[
 \Theta_n=\pi\frac{F_{n-1}}{F_{n+1}}
 \longrightarrow\frac{\pi}{\varphi^2}.
\]
La convergence se déduit de la valeur propre dominante \(\varphi\), ou des
quotients consécutifs déjà bornés. Le lecteur
\(\mathcal R(x,y)=x/y^2\), sous l’échange des arguments,
exprime également
\begin{equation}
 \rho_A=\frac{\pi}{\varphi^2},\qquad
 \rho_E=\frac{\varphi}{\pi^2},\qquad
 \varphi^3\rho_A^2\rho_E=1.
 \label{vi:dep:dos-orientaciones}
\end{equation}
La dernière égalité se vérifie en multipliant les fractions.
Elle relie deux orientations du même lecteur et ne sélectionne pas les
autres coefficients de l’équation électronique.

Le retour central typé a pour longueur \(27\), avec un support
\(\ZZ/9\ZZ\times\FF_3\). Les cinq positions de la microfibre
s’insèrent, dans la jauge régionale fixée, comme
\((4,0),(5,0),(6,0),(6,1),(6,2)\).
Le complément possède \(27-5=22\) positions. Son incidence avec les
trois états de phase tritique comporte \(5\cdot3=15\) paires.
On obtient ainsi les entiers \(-22\) et \(15\) utilisés par le
lecteur électronique ; l’orientation fixe le signe du premier terme.
La composition cubique utilise trois copies de la coordonnée de fermeture,
de sorte que son évaluation est \(\alpha^3\). Le support, l’insertion
régionale et cette règle de composition sont des données opératoires distinctes ;
un dénombrement isolé ne remplace pas leur déclaration.

Le défaut global et son retour nonadique s’expriment par
\begin{equation}
 d_4=\frac{\pi-e\log\pi}{270},\qquad
 \Delta_4=d_4\left(1+\frac{\pi}{729}\right).
 \label{vi:dep:delta4}
\end{equation}
Le dénominateur \(270\) appartient à la couronne privée de son unité de
clôture ; \(729=9^3\) au support cubique. La fonction
\(f(x)=x-e\log x\) atteint son minimum nul en \(x=e\), car
\(f'(x)=1-e/x\) y change de signe. Comme \(\pi\ne e\),
on obtient \(d_4>0\) et \(\Delta_4>0\).

\begin{definition}[Coordonnée électronique basale]
\label{vi:dep:electron-basal}
La composition des lecteurs précédents est
\begin{equation}
 \varepsilon_e^{(0)}=
 \frac{\sqrt3}{4}
 \left[\exp\left(\frac{\varphi}{\pi^2}\right)-22\alpha^3\right]
 (1+15\Delta_4).
 \label{vi:dep:formula-basal}
\end{equation}
\end{definition}
L’exponentielle agit sur l’orientation \(\rho_E\) de
\eqref{vi:dep:dos-orientaciones} ; la lecture de vacance corrige cette
composante et le dernier facteur conserve le retour du défaut global.
Pour \(0<\alpha<0.01\), tous les facteurs sont positifs :
\(\exp(\varphi/\pi^2)>1\), \(22\alpha^3<22\cdot10^{-6}<1\)
et \(\Delta_4>0\). C’est une coordonnée adimensionnelle. Sa présence
dans une coordinatisation de masse est définie après la construction de l’échelle.

\subsection{Échelle déterminantale et constante de Planck réduite}
\label{vi:dep:planck}

Le bloc local \eqref{vi:dep:hadamard} a pour déterminant \(-16\).
Sa forme normale de Smith est \(\operatorname{diag}(1,2,2,4)\).
En effet, le plus grand diviseur des entrées est un ; les mineurs
d’ordre deux sont pairs et l’un vaut deux en valeur absolue. Les mineurs
d’ordre trois sont, par l’identité de la comatrice, de valeur absolue quatre,
et le déterminant a pour valeur absolue seize. Les diviseurs
déterminantaux sont \(1,2,4,16\), dont les rapports successifs donnent les
quatre facteurs de Smith. Par conséquent,
\begin{equation}
 G_H=\operatorname{coker}H_4
 \simeq\ZZ/2\ZZ\oplus\ZZ/2\ZZ\oplus\ZZ/4\ZZ,
 \qquad |G_H|=16.
 \label{vi:dep:conucleo}
\end{equation}

Le lecteur normé local multiplie la même coordonnée \(\alpha\)
une fois pour chaque classe de ce conoyau et applique une fois la coordonnée
d’auto-échelle. Sa définition produit
\begin{equation}
 \Sigma_H=\varphi\prod_{g\in G_H}\alpha
         =\varphi\alpha^{16}.
 \label{vi:dep:sigma-H}
\end{equation}
L’ordre du conoyau détermine l’exposant de ce lecteur. La forme
normale de Smith et le choix du lecteur normé remplissent des fonctions
distinctes : la première ne transforme pas n’importe quelle fonctionnelle du bloc en
l’expression \eqref{vi:dep:sigma-H}.

\begin{proposition}[Décade d’action]
\label{vi:dep:decada}
La sortie \eqref{vi:dep:sigma-H} satisfait
\(10^{-34}<\Sigma_H<10^{-33}\) ; sa valuation décimale est \(-34\).
\end{proposition}
\begin{proof}
Les évaluations régionales et positionnelles donnent les intervalles
\(1.618<\varphi<1.619\) et \(0.007297<\alpha<0.007298\).
Par positivité,
\[
 \frac{1618\,7297^{16}}{10^{99}}
 <\Sigma_H<
 \frac{1619\,7298^{16}}{10^{99}}.
\]
Les inégalités entières
\(1618\,7297^{16}>10^{65}\) et
\(1619\,7298^{16}<10^{66}\) prouvent le résultat sans utiliser de
constante d’action externe.
\end{proof}

La décade et la mantisse sont des représentations d’étapes différentes.
Avant de construire la seconde, le lecteur régional sélectionne le retour
persistant de la microfibre de clôture. La projection
\(p_{\rm ret}(u_1,\ldots,u_6)=(u_4,u_5,u_6)\) appliquée aux
cinq émissions de \ref{vi:dep:regiones} produit le bloc
\(b_\pi=(169,272,272)\) quatre fois et
\(b_\pi^-=(418,674,521)\) une fois. Ces multiplicités dénombrent des
lignes régionales, non les multiplicités des germes \(432,144\).

\begin{proposition}[Sélecteur équivariant du retour régional]
\label{vi:dep:selector169}
Parmi les règles qui conservent les permutations des cinq régions,
sélectionnent le bloc terminal de multiplicité maximale, sont équivariantes
sous les permutations de ses trois positions et retiennent une position fixe
sous le stabilisateur du bloc, la coordonnée sélectionnée est unique
et vaut \(\rho_\pi=169\).
\end{proposition}
\begin{proof}
La projection des cinq lignes montre que le bloc de multiplicité
maximale est unique et est \(b_\pi\). Son stabilisateur dans
\(\mathfrak S_3\) échange les deux positions de valeur \(272\)
et fixe uniquement la position de valeur \(169\). Retenir cette position
est compatible avec toute permutation simultanée des coordonnées.
De manière équivalente, la valeur s’exprime sans choisir de position initiale
comme
\(\operatorname{sing}(b_\pi)=\sum_j(b_\pi)_j
-2\operatorname{moda}(b_\pi)=169\).
\end{proof}

La graduation par longueur exprime un mot de longueur \(n\)
au moyen du caractère \(\chi_z([w])=z^{|w|}\). La concaténation
additionne les longueurs, de sorte que
\(\chi_z([uv])=\chi_z([u])\chi_z([v])\). Comme les émissions
régionales appartiennent à \(\mathcal U_6\), l’impulsion du retour
apparaît comme \(\rho_\pi z^6\). Le degré six correspond à la
longueur de l’émission, non à la longueur d’une marche fermée qui
parcourt plusieurs ancrages du graphe.

Le poids de prolongement s’obtient dans la droite déterminant du
transport bicouche. Pour une involution réelle \(R\) en dimension deux,
avec \(\operatorname{tr}R=0\), soient
\[
 x_A=\frac{50\pi}{9}\alpha,\qquad
 \mathcal T(x_A,y)=\exp(-x_AI_2-yR).
\]
La coordonnée orientée \(y\) est encore formelle. Dans une base propre
de \(R\), les deux entrées sont \(e^{-x_A-y}\) et
\(e^{-x_A+y}\). Par conséquent,
\[
 \bigwedge^2\mathcal T=\det\mathcal T
 =e^{-2x_A}=D_A.
\]
La dépendance en \(y\) s’annule. La multiplicativité par
concaténation détermine le caractère de prolongement \(k\mapsto D_A^k\) :
tout caractère de valeur élémentaire \(D_A\) doit avoir cette valeur
en \(k=1+\cdots+1\). Cette construction précède l’évaluation
de la coordonnée angulaire conjuguée.

Le lecteur régional d’ordre cinq emploie les valeurs propres \(9,5\)
du bloc central \(\bigl(\begin{smallmatrix}7&2\\2&7\end{smallmatrix}\bigr)\),
l’orbite de quatre positions et le calendrier de douze fenêtres.
Sa règle graduée fixe
\begin{align}
 H_5={}&\frac12\sqrt{\frac{1000\alpha}{\varphi}}-\alpha
       +\frac9{16}\alpha^2-\frac59\alpha^3
       +\frac7{48}\alpha^4-\frac1{54}\alpha^5,\notag\\
 D_A={}&\exp\left(-\frac{100\pi\alpha}{9}\right),\notag\\
 \mathcal C_\pi={}&169\alpha^6+
                  \frac{D_A\alpha^7}{1-D_A\alpha},\notag\\
 \eta_{\rm ret}={}&H_5-\frac{90}{\pi}\mathcal C_\pi.
 \label{vi:dep:mantisas}
\end{align}
Dans cette règle, les quotients sont
\(9/16=9/4^2\), \(5/9\), \(7/48=(\operatorname{tr}B_c/2)/(4\cdot12)\)
et \(1/54=1944/104976\). Leur affectation aux degrés et les signes
font partie du lecteur gradué. L’entier \(169\) est la représentation
du sélecteur démontré dans la proposition \ref{vi:dep:selector169}.
Sa donnée régionale est antérieure à l’évaluation de \(\alpha\) ; le
lecteur analytique la compose ensuite avec la coordonnée de fermeture.

La queue de \(\mathcal C_\pi\) possède une construction par renouvellement.
L’impulsion \(\rho_\pi z^6\) et la continuation stricte appartiennent
à des termes distincts. Après avoir enregistré une seule fois l’impulsion,
la continuation est normalisée par l’unité de la droite déterminant ;
chaque nouveau pas la multiplie par \(D_A\). L’orientation de la
jauge régionale place ce retour dans le secteur compensateur qui
est soustrait du caractère de cinquième degré.
Avec un paramètre formel \(z\), la règle
\begin{equation}
 R(z)=D_Az^7+D_AzR(z)
 \label{vi:dep:renovacion}
\end{equation}
détermine de manière unique
\(R(z)=D_Az^7/(1-D_Az)\). Cela peut être prouvé par comparaison des
coefficients ou par inversion de la série formelle \(1-D_Az\).
La normalisation du premier retour, outre le rapport récurrent,
est essentielle à cette unicité. Comme \(0<D_A<1\) et \(\alpha<1\),
l’évaluation en \(z=\alpha\) converge absolument.

\begin{proposition}[Positivité et rapport de retour d’action]
\label{vi:dep:positividad-accion}
Sur la branche engendrée, \(0<\eta_{\rm ret}<H_5\). Le rapport
\begin{equation}
 R_{\rm act}=\frac{H_5}{\eta_{\rm ret}}
 \label{vi:dep:Ract}
\end{equation}
est positif, supérieur à un et adimensionnel.
\end{proposition}
\begin{proof}
Les bornes \(0.007<\alpha<0.008\) et \(\varphi<1.619\) donnent
\(\tfrac12\sqrt{1000\alpha/\varphi}>1\). La somme des valeurs absolues
des termes négatifs restants de \(H_5\) est inférieure à
\(0.008001\), de sorte que \(H_5>0.99\).
Comme \(D_A<1\), \(\pi>3\) et \(\alpha<0.008\),
\[
 0<\frac{90}{\pi}\mathcal C_\pi
 <30\left(169(0.008)^6+
                  \frac{(0.008)^7}{1-0.008}\right)<10^{-6}.
\]
On en déduit \(0<\eta_{\rm ret}<H_5\). Le rapport de deux coordonnées
d’action dans la même base est adimensionnel et ne change pas lors d’un changement d’échelle
de cette base.
\end{proof}

Soit \(\mathcal U_S\) une base de la droite d’action. Les deux
sections, antérieure et postérieure au retour, sont
\begin{equation}
 \hbar_{\rm pre}=H_5\,10^{-34}\mathcal U_S,
 \qquad
 \hbar_{\rm ret}=\eta_{\rm ret}\,10^{-34}\mathcal U_S.
 \label{vi:dep:hbar}
\end{equation}
Sa lecture physique est la construction proposée de la constante de Planck
réduite, avec la section de retour explicitée. L’unité
\(\mathcal U_S\mapsto\mathrm{J\,s}\) est une réalisation métrologique
postérieure. Le coefficient de \(\Sigma_H\) dans sa décade n’est pas
\(H_5\) : le déterminant exprime l’échelle et le lecteur régional
exprime la coordonnée qui occupe cette échelle. Pour chaque section,
\(h=2\pi\hbar\) exprime l’action d’un tour complet.

\subsection{Coordonnées angulaires et horloge de la réalisation absolue}
\label{vi:dep:reloj}

La paire angulaire utilise les coordonnées adimensionnelles déjà engendrées :
\begin{equation}
 A_{\deg}=1000\alpha,\qquad
 C^*_{\deg}=2(\eta_{\rm ret}+\alpha),\qquad
 x=\frac{\pi}{180}A_{\deg},\quad
 y=\frac{\pi}{180}C^*_{\deg}.
 \label{vi:dep:angulos}
\end{equation}
La première opération décale d’un bloc le développement en base mille.
La seconde est la représentation bidirectionnelle de la coordonnée régionale :
de manière équivalente,
\(y=\pi(H_5+\alpha)/90-\mathcal C_\pi\).
Cette égalité se déduit en substituant \eqref{vi:dep:mantisas} ; elle
n’introduit pas de règle supplémentaire de sélection du contre-angle.
La coordinatisation d’un tour comporte \(360\) unités angulaires ou \(2\pi\)
radians. Ainsi, pour une coordonnée relevée,
\[
 \hbar\theta_{\rm rad}
 =h\frac{\theta_{\deg}}{360}.
\]
Le facteur \(10^{-34}\) reste dans la droite d’action et ne
s’insère pas dans les coordonnées angulaires.

Les canaux positifs sont
\[
 q_+=e^{-x-y},\qquad q_-=e^{-x+y}.
\]
On récupère la paire par
\(x=-\tfrac12\log(q_+q_-)\),
\(y=\tfrac12\log(q_-/q_+)\). L’échange des canaux conserve
\(x\) et change le signe de \(y\). L’opérateur muni des projecteurs
orientés \(P_\pm\),
\(T=q_+P_++q_-P_-\) conserve ces étiquettes en exprimant ses deux
valeurs propres.

Pour une unité temporelle interne \(t_0>0\), la fermeture de \(108\)
pas détermine
\begin{equation}
 \omega_0=\frac{2\pi}{108t_0},\qquad
 E_{\rm clk}=\hbar_{\rm ret}\omega_0,
 \qquad m_{\rm clk}=\frac{E_{\rm clk}}{c_{\rm int}^{2}}.
 \label{vi:dep:reloj-absoluto}
\end{equation}
Ces formules produisent des dimensions de fréquence, d’énergie et de masse.
La vitesse et la longueur s’obtiennent par la réponse des
mêmes canaux, comme le précisent les identités suivantes. Les exprimer
en secondes, en joules ou en unités de masse nécessite les coordinatisations
correspondantes ; ces coordinatisations ne modifient ni les canaux ni la section
électronique qui est normalisée par rapport à l’horloge.

\subsubsection{Identité entre vitesse constitutive et vitesse de l’horloge}

Dans le domaine \(x>|y|\), les deux canaux satisfont \(0<q_\pm<1\).
Sur la réalisation à deux feuillets, munie de projecteurs de rang un
\(P_+,P_-\), on définit la réponse temporelle
\[
 \mathsf V=(I-T^{90})(I-T^{120})^{-1}
          =r_+P_++r_-P_-,
 \qquad r_\pm=\frac{1-q_\pm^{90}}{1-q_\pm^{120}}>0.
\]
L’inverse existe car les deux valeurs propres de \(T^{120}\) sont inférieures
à un. La lecture symétrique de la réponse est
\[
 \mathcal E=
 \log[(1-q_+^{90})(1-q_-^{90})]
 -\log[(1-q_+^{120})(1-q_-^{120})]
 =\log(r_+r_-).
\]
Soient \(u_S,u_C,u_T\) les bases positives d’action, de vitesse et de
temps, et \(u_L=u_Cu_T\) la base de longueur. Les lectures
cinématique et constitutive sont
\[
 c_{\rm int}=e^{-\mathcal E}u_C,\qquad
 c_{\rm vac}=(r_+r_-)^{-1}u_C.
\]

\begin{proposition}[Réalisation commune de propagation et de masse absolue]
\label{vi:dep:velocidad-comun}
Dans une même coordinatisation dimensionnelle,
\[
 c_{\rm int}=c_{\rm vac}=\widehat c\,u_C,
 \qquad \widehat c=(\det\mathsf V)^{-1}=(r_+r_-)^{-1}.
\]
Pour \(t_0=\tau_0u_T\), avec \(\tau_0>0\), la longueur élémentaire est
\[
 \ell_0=c_{\rm int}t_0=\tau_0\widehat c\,u_L.
\]
La masse chronologique de \eqref{vi:dep:reloj-absoluto} conserve, par
conséquent, la forme
\[
 m_{\rm clk}=
 \frac{2\pi\,\eta_{\rm ret}\,10^{-34}}
      {108\,\tau_0\,\widehat c^{\,2}}\,
 \frac{\mathcal U_S}{u_Tu_C^2}.
\]
\end{proposition}
\begin{proof}
L’identité \(\mathcal E=\log(r_+r_-)\) donne
\(e^{-\mathcal E}=(r_+r_-)^{-1}\) ; le rang un de chaque feuillet donne
\(\det\mathsf V=r_+r_-\). L’échange des feuillets permute les
facteurs et conserve leur produit. La longueur s’obtient en multipliant
la même vitesse par \(t_0\). Substituer ces expressions et
\eqref{vi:dep:hbar} dans \eqref{vi:dep:reloj-absoluto} produit la
formule massique. Le facteur \(\mathcal U_S/(u_Tu_C^2)\) a la
dimension d’une masse ; il n’intervient pas dans la sélection des canaux.
\end{proof}

Pour une route homogène de \(n>0\) transitions, les lectures additives
\(T(\gamma)=nt_0\) et \(L(\gamma)=n\ell_0\) donnent
\(L(\gamma)/T(\gamma)=c_{\rm vac}\). Si le canal varie selon
l’arête, on additionne leurs longueurs ; le quotient est alors la vitesse
moyenne de ces arêtes, pas nécessairement une vitesse locale commune.

L’égalité des sections demeure sous les changements de coordinatisation. Si
\(u_C'=du_C\), \(u_T'=\eta u_T\), avec \(d,\eta>0\), alors
\[
 \widehat c'=\widehat c/d,\qquad
 \tau_0'=\tau_0/\eta,\qquad u_L'=d\eta u_L,
 \qquad
 \tau_0'\widehat c'u_L'=\ell_0.
\]
En particulier, choisir la base adaptée \(u_C'=c_{\rm vac}\)
rend \(\widehat c'=1\), mais la coordonnée du lecteur dans la
coordinatisation interne reste \((r_+r_-)^{-1}\).

\subsection{Deux lecteurs de la trajectoire électronique et évaluation bêta}
\label{vi:dep:beta}

La trajectoire centrale commence en \((r,c)=(5,5)\), avec
l’orientation \((h,v)=(1,1)\). TRIT répète le cycle \((+1,-1,0)\).
Dans le premier régime, on lit \(r+c\) puis on avance de
\(c\mapsto1+((c-1+h)\bmod9)\) ; dans le second, on lit \(rc\) puis
on avance de \(r\mapsto1+((r-1+v)\bmod9)\) ; dans le régime
zéro, on lit \(9\), sans déplacement. On exprime la classe modulo
trois de la racine numérique de la lecture. Après les pas \(27,54,81\),
on inverse simultanément \(h,v\). La représentation de \(108\)
pas conserve la trajectoire complète et donne
\begin{equation}
 a=100000220,\quad b=120200000,\quad
 \gamma_e=(a^3b^3)^2,
 \quad\tau_e=(+,+,+,-,-,-,+,+,+,-,-,-).
 \label{vi:dep:trayectoria-electron}
\end{equation}
L’identité se vérifie en parcourant les neuf instants de chaque
fenêtre ; au troisième bloc, l’inversion remplace \(a\) par \(b\),
et le second parcours de \(54\) pas répète la représentation.
L’état conserve également ses coordonnées de mémoire.

Le premier lecteur mesure les désaccords cycliques de la représentation ternaire :
\[
 D_k=\#\{t\in\ZZ/108\ZZ:\gamma_e(t)\ne\gamma_e(t+k)\}.
\]
Comme le mot exprimé est de période \(54\), chaque désaccord
apparaît avec son translaté de \(54\) ; tous les \(D_k\) sont pairs.
Dans la moitié \(a^3b^3\), le dénombrement des désaccords pour
\(k=1,3,4,27,36\) donne respectivement \(27,28,34,24,16\).
En le doublant,
\[
 (D_1,D_3,D_4,D_{27},D_{36})=(54,56,68,48,32).
\]
Par conséquent,
\begin{equation}
 K_D=\left(D_{27}+D_{36},D_1,\frac{D_4-D_3}{2}\right)
     =(80,54,6).
 \label{vi:dep:KD}
\end{equation}

Le second lecteur agit sur l’orientation des douze fenêtres.
Définissons, avec des indices cycliques,
\[
 C_k=\sum_{r=0}^{11}\tau_r\tau_{r+k},\quad
 A_\ell=\sum_{r=0}^{11}\left(\sum_{j=0}^{\ell-1}\tau_{r+j}\right)^2,
 \quad P_6=\sum_{r=0}^{5}\tau_r\tau_{r+6}.
\]
La combinaison entière primitive orientée qui annule la partie diagonale
de \(A_3,A_4\) est \(\Omega=4A_3-3A_4\) : ses coefficients
satisfont \(3u+4v=0\), \(\gcd(u,v)=1\), et l’on fixe \(u>0\).
Le développement du carré prouve
\[
 A_\ell=\ell C_0+
     2\sum_{k=1}^{\ell-1}(\ell-k)C_k,
 \qquad\Omega=-2C_1-4C_2-6C_3.
\]
Pour \(\tau_e\), on obtient
\(C_0=12,C_1=4,C_2=-4,C_3=-12\),
\(A_3=44,A_4=32,P_6=6\). D’où
\begin{equation}
 K_\Omega=(\Omega,54,P_6)=(80,54,6)=K_D.
 \label{vi:dep:KOmega}
\end{equation}
La seconde composante conserve le demi-retour de \(108\) pas ;
elle n’est pas introduite par l’annulation diagonale. L’égalité des deux
lecteurs a été prouvée sur la trajectoire centrale spécifiée ;
elle ne s’étend pas par définition à toutes les collections de l’atlas.

\begin{definition}[Coordonnée électronique postérieure au retour]
La composition électronique complète est
\begin{equation}
 \varepsilon_e^\beta
  =\varepsilon_e^{(0)}
      R_{\rm act}^{80-54\alpha+6\Delta_4}.
 \label{vi:dep:electron-beta}
\end{equation}
\end{definition}
La positivité de \eqref{vi:dep:formula-basal} et
\eqref{vi:dep:Ract} permet de définir cette puissance au moyen du
logarithme réel. Le triplet entier provient indifféremment des deux
lecteurs centraux qui viennent d’être calculés. L’évaluation relative
préserve la constante de Planck réduite par le rapport de ses
sections, et la réalisation absolue utilise en outre
\eqref{vi:dep:reloj-absoluto} et la coordinatisation dimensionnelle déclarée.

Dans les tableaux hérités, \(\mathfrak e_0\) désigne la
coordonnée basale \(\varepsilon_e^{(0)}\) dans la coordinatisation qui y est
spécifiée. La section \(\hbar_{\rm int}\) utilisée pour une
évaluation de survie doit être déclarée dans la même coordinatisation :
pour la réalisation postérieure au retour du présent développement,
c’est \(\hbar_{\rm ret}\). Les lignes historiques conservent la
convention déclarée dans leur tableau d’origine et ne sont pas recalculées par un
changement silencieux de cette section.

Si \(\mathcal U_E\) est la base d’énergie de l’évaluation
électronique, \(E_e=\varepsilon_e^\beta\mathcal U_E\). Lorsqu’on emploie
le quantum chronologique comme base, le coefficient est
\[
 \widehat\varepsilon_e
   =\varepsilon_e^\beta\frac{\mathcal U_E}{E_{\rm clk}},
 \qquad E_e=\widehat\varepsilon_e E_{\rm clk}.
\]
L’égalité conserve l’objet et transforme sa coordonnée. Ainsi,
le nombre présenté dans une coordinatisation en MeV ne peut pas être multiplié
sans conversion par n’importe quel autre quantum d’énergie. L’électron
conserve sa construction structurelle lorsqu’il est exprimé par rapport à une
échelle chronologique : il n’est pas remplacé par la masse de ce quantum.

La comparaison métrologique s’applique au résultat de cette composition,
avec la version de \(\alpha\), la section d’action, le lecteur
\(\Delta_4\) et la coordinatisation des unités déclarées. En particulier,
\(d_4\) et \(\Delta_4\) ne sont pas interchangeables : les remplacer
simultanément dans le facteur basal et dans l’exposant modifie le
résultat selon le rapport positif
\[
 \frac{1+15\Delta_4}{1+15d_4}
 R_{\rm act}^{6(\Delta_4-d_4)}>1.
\]
Les formules conservées permettent de localiser cette différence sans
l’attribuer à une masse externe ni sélectionner à nouveau la trajectoire.
